\documentclass[11pt]{article}
\usepackage[utf8]{inputenc}
\usepackage{amsmath,amssymb}
\usepackage{graphicx}
\usepackage{booktabs}
\usepackage{authblk}
\usepackage[margin=1in]{geometry}
\usepackage[hidelinks]{hyperref}

\newcommand{\TODO}[1]{\textbf{[TO MEASURE: #1]}}

\title{A 3D Multiphase Continuum Computational Model for Atoms}

\author[1]{Claes Johnson}
\affil[1]{KTH Royal Institute of Technology, Stockholm, Sweden \\
\texttt{claesjohnson@gmail.com}}
\date{\today \\[0.4em]
\small\textit{Mathematical formulation by the author; implementation developed} \\
\small\textit{with AI assistance (Claude, Anthropic).}}

\begin{document}
\maketitle

\begin{abstract}
RealQM is a general, in-principle parameter-free, low-cost computational method for
quantum mechanics, formulated as three-dimensional continuum mechanics. The electrons
form a sum of one-electron components, each carrying unit charge on a non-overlapping
domain separated by free (Bernoulli) boundaries, and the ground state minimises the
standard Coulomb energy over both the components and the domain partition. The method is
applied here to the periodic table, computing total energies and first ionization
energies of all atoms in spherical symmetry, with the whole table covered in about a
minute on a laptop. Each electron carries its own energy, so the first ionization energy
is read directly as the valence-electron energy, without a separate ion calculation.
Total energies fall within about one to two percent across the table and ionization
energies within ten to twenty percent, with the correct ionization ordering throughout,
including 4s before 3d across the transition series.
\end{abstract}

\medskip
\noindent\textbf{Keywords:} atomic structure; real-space methods; multiphase wave
functions; free-boundary problems; ionization energies; imaginary-time propagation;
electron packing.

\section{Introduction}\label{sec:intro}

This paper applies \textbf{RealQM} --- a 3D continuum-mechanical model of the atom in
which the $N$ electrons are represented not as a single antisymmetric wave function on
$3N$-dimensional configuration space but as a sum of $N$ one-electron components on
mutually non-overlapping domains of physical three-dimensional space, separated by a free
boundary, both the components and the partition being unknowns of a single Coulomb-energy
minimisation --- \emph{in spherical symmetry}, as a detailed test on the \textbf{periodic
table}: the total atomic energies, electron packing, and ionization energies of the atoms
across all periods, taken from the outset as its own subject. The model and its wider
development --- molecules, bonding, atomization energies, reactions --- are given in the
de Broglie article~\cite{Johnson-RealQM} and the accompanying monograph; the present
paper takes the \emph{atoms} themselves as its subject, in periodic-table detail.

\paragraph{Context: one method across the whole table.} High-accuracy atomic reference
data are assembled from a family of specialised treatments, each chosen to suit the
system at hand: correlated non-relativistic wave-function methods for the light
atoms~\cite{SzaboOstlund,ClementiRoetti,Chakravorty}, relativistic Dirac--Fock for the
heavy ones~\cite{Desclaux}, density-functional and real-space approaches for large-scale
work~\cite{HK64,KS65,Beck2000}, alongside tabulated experimental values~\cite{NIST-ASD};
these reach accuracies far beyond the present model. The test reported here is of a
different kind. A single, in-principle parameter-free scheme is applied uniformly to
\emph{every} atom --- solved the same way for hydrogen and for uranium, with no
per-element adjustment and no fitted functional --- and asked to reproduce the gross
electronic structure of the whole table at once. What a uniform test of this sort can
establish is not competitive accuracy on any one atom but \emph{structural correctness
across scope}: that one Coulomb-energy minimisation recovers the shell filling, the
total-energy trend to a few percent, and --- the more discriminating outcome --- the
ionization-energy periodicity and ordering, including $4s$ before $3d$ across the
transition series. Uniformity across the table, rather than precision on a single atom,
is the property put to the test.

\paragraph{Ionization energy as a per-electron quantity.} The feature that makes
atoms a natural target is that, in this multiphase formulation, each electron
satisfies its own one-electron eigenproblem and therefore carries its own energy
$E_i$ (Section~\ref{sec:model}). The ionization energy is then the energy of the
valence electron \emph{directly},
\begin{equation}\label{eq:IE-direct}
\mathrm{IE} \;=\; |E_{\rm valence}|,
\end{equation}
not a difference $E(\mathrm{A}^+)-E(\mathrm{A})$ of two large total energies. This
is consequential numerically. An atomic total energy is large (e.g.\ Ne:
$\approx -128.9$ Ha) while an ionization energy is small ($\mathrm{IE}_1$ of Ne
$\approx 0.79$ Ha); computing the latter as a difference of the former demands the
total energy to absolute accuracy far beyond $1\%$. The per-electron route avoids
this entirely: the spectator core never enters~\eqref{eq:IE-direct}, so ionization
energies can be accurate even when absolute total energies are only at the
${\sim}1\%$ level. This division --- loose total energies, sharp ionization
energies --- is the organising idea of the paper.

\paragraph{Contributions.} (i) A statement of the multiphase 3D atom model and of
the per-electron energy $E_i$ as the computed ionization energy
(Section~\ref{sec:model}); (ii) a two-scale solution architecture pairing a
spherical inner-shell solve with a valence solve in the resulting pseudo-kernel ---
spherical for all results here, full 3D~\cite{Johnson-RealQM} for the two geometries
that require it (Section~\ref{sec:twoscale}); (iii) the solution algorithm
(Section~\ref{sec:algorithm}); and (iv) validation
across the periodic table --- first and second ionization energies against NIST,
and total energies at the ${\sim}1\%$ level --- together with the numerical
properties of the solver (Sections~\ref{sec:results}--\ref{sec:convergence}).

\section{The multiphase atom model}\label{sec:model}

We summarise only what is needed here; derivations are
in~\cite{Johnson-RealQM}. For an atom of nuclear charge $Z$ the $N$-electron wave
function is
\begin{equation}\label{eq:Psi-sum}
\Psi(\mathbf{x}) = \sum_{i=1}^{N} \psi_i(\mathbf{x}),
\qquad \psi_i \in H^1(\Omega_i),
\qquad \int_{\Omega_i} \psi_i^2 \, d\mathbf{x} = 1,
\end{equation}
on a partition of $\mathbb{R}^3$ into non-overlapping domains $\Omega_i$ with free
boundaries $\Gamma_i=\partial\Omega_i$. The total energy is the standard Coulomb
functional
\begin{equation}\label{eq:TE}
TE(\Psi)=\underbrace{\tfrac12\sum_i\int_{\Omega_i}|\nabla\psi_i|^2}_{\text{kinetic}}
\;+\;\underbrace{\int_{\mathbb{R}^3} W\,\Psi^2}_{\text{kernel--electron}}
\;+\;\underbrace{\sum_i\int_{\Omega_i}\sum_{j\ne i}V_j\,\psi_i^2}_{\text{electron--electron}},
\end{equation}
with kernel potential $W(\mathbf{x})=-Z/|\mathbf{x}|$ and per-electron Hartree
potentials
$V_j(\mathbf{x})=\tfrac12\int_{\Omega_j}\psi_j^2(\mathbf{y})/|\mathbf{x}-\mathbf{y}|\,d\mathbf{y}$;
self-repulsion is excluded by construction ($j\ne i$ on disjoint domains).
Stationarity in $\psi_i$ at fixed $\Omega_i$ gives the one-electron eigenproblem
\begin{equation}\label{eq:Hi}
\mathcal{H}_i\psi_i\equiv\Bigl(-\tfrac12\Delta + W + 2\!\sum_{j\ne i}V_j\Bigr)\psi_i = E_i\psi_i
\ \text{ in }\Omega_i,
\qquad \frac{\partial\psi_i}{\partial n}=0 \ \text{ on }\Gamma_i,
\end{equation}
with continuity of $\psi_i$ across the inter-electron interface and the homogeneous
Neumann condition together constituting the \textbf{Bernoulli free-boundary
condition}. Spatial exclusion plays the role of the Pauli principle: the electrons
occupy disjoint territories rather than antisymmetrised orbitals.

\paragraph{The per-electron energy is the ionization energy.} Equation~\eqref{eq:Hi}
gives each electron a well-defined energy $E_i$ --- its kinetic energy plus its
attraction to the kernel plus its interaction with the field of the other
electrons. For the valence electron, $|E_{\rm valence}|$ is the energy required to
remove it, i.e.\ the ionization energy~\eqref{eq:IE-direct}. The second ionization
energy is obtained from the valence energy of the singly-ionized atom, recomputed
with the outer electron removed and the remaining partition relaxed. Because the
inner electrons are spectators to this process, the accuracy of $\mathrm{IE}$ is
controlled by the valence solve and not by the (looser) accuracy of the core.

\section{Two-scale solution: spherical core, 3D valence}\label{sec:twoscale}

The per-electron structure permits a computation that mirrors the physics: the
spectator core is treated cheaply in spherical symmetry, and only the valence
electrons --- which determine the ionization energies --- are treated explicitly in
3D.

\paragraph{Inner shells (spherical).} The closed inner shells are solved by a
one-dimensional radial reduction of \eqref{eq:Psi-sum}--\eqref{eq:Hi}: radial
densities $u(r)$ with free boundaries between shells set by continuity of $u$. The
converged core supplies an effective \textbf{pseudo-kernel}
\begin{equation}\label{eq:pseudokernel}
V_{\rm kernel}(r) = -\frac{Z_{\rm kernel}}{r}\ \ (r>r_c),
\qquad Z_{\rm kernel} = Z - N_{\rm core},
\end{equation}
smoothly softened for $r\le r_c$, where $Z_{\rm kernel}$ is the nuclear charge
screened by the $N_{\rm core}$ core electrons and the softening radius $r_c$ is read
off the converged core density. Crucially, $r_c$ is thereby \emph{determined by the
core solve rather than fitted}, so the valence calculation carries no per-element
empirical parameter.

\paragraph{Valence electrons.} In the general two-scale method the valence electrons
are solved explicitly on a 3D grid in the pseudo-kernel~\eqref{eq:pseudokernel},
retaining the angular (sub-shell) structure the radial reduction cannot represent.
\emph{In this paper the valence is solved in the same spherical reduction}
(Section~\ref{sec:radialsolver}); the full 3D solve is reserved for the two geometries
that require it (the closed octet and the interpenetrating $4s/3d$). Either way the
per-electron energies $E_i$ from~\eqref{eq:Hi} give the ionization energies
through~\eqref{eq:IE-direct}, and the loose ${\sim}1\%$ accuracy of the spherical core
enters only through the effective potential it presents to the valence, to which
$\mathrm{IE}$ is far less sensitive than to the core energy itself.

\subsection{The spherically-reduced solver}\label{sec:radialsolver}

Every result in this paper is computed with the spherical (1D radial) reduction of the
model, which we state explicitly. The atom has shells $s=0,\dots,S{-}1$; shell $s$
carries $N_s$ electrons as a radial amplitude $u_s(r)$ on $[R_{s-1},R_s]$ (with
$R_{-1}=0$, $R_{S-1}=\infty$). With the $s$-wave radial Laplacian
$\Delta u = u'' + \tfrac{2}{r}u'$ and Hartree atomic units, the per-electron operator,
its imaginary-time relaxation to the ground state, and the per-electron Hartree
potential of the \emph{other} shells are
\begin{align}
H_s &= -\tfrac12\Delta - \frac{Z}{r} + W_s(r), &
\frac{\partial u_s}{\partial\tau} &= -H_s u_s, \label{eq:radial-itp}\\
-\Delta W_s &= 4\pi\!\!\sum_{t\neq s}\! |u_t|^2, &
&W_s(r)=\tfrac1r\!\int_0^r\!\!\rho_{\neq s}4\pi r'^2dr' + \!\int_r^\infty\!\!\rho_{\neq s}4\pi r'dr', \label{eq:radial-hartree}
\end{align}
the same-shell self-term reduced by $(N_s{-}1)/N_s$. The interfaces $R_s$ are free
(Bernoulli) boundaries, fixed by continuity of the amplitude, and each shell is
normalised:
\begin{equation}\label{eq:radial-bc}
u_s(R_s)=u_{s+1}(R_s), \qquad 4\pi\!\int_0^\infty |u_s|^2 r^2\,dr = N_s .
\end{equation}
The total energy (half-Hartree, no double counting) and the per-electron eigenvalue
(full Hartree) --- the latter the ionization energy of Section~\ref{sec:model} --- are
\begin{equation}\label{eq:radial-E}
E=\sum_s 4\pi\!\int\!\Big[\tfrac12|u_s'|^2-\tfrac{Z}{r}|u_s|^2+\tfrac12 W_s|u_s|^2\Big]r^2dr,
\qquad \varepsilon_s=\frac{\langle u_s|H_s|u_s\rangle}{N_s},\ \ \mathrm{IE}=|\varepsilon_{\rm valence}|.
\end{equation}
The scheme is discretised on a uniform grid $r_i=ih$ with an explicit step
$\Delta\tau=\tfrac12 h^2$ (first order, cusp-limited; Section~\ref{sec:convergence}).
For $Z>18$ the nuclear term is replaced by the pseudo-kernel
$Z/r\to Z_{\rm ker}/\max(r,r_c)$ of~\eqref{eq:pseudokernel} with a Neumann inner
boundary $u'(r_c)=0$, confining the valence outside the core.

\section{Algorithm: the three-line solver}\label{sec:algorithm}

The radial system of Section~\ref{sec:radialsolver} is solved by three coupled
relaxations to a joint steady state --- the \emph{three-line solver}. At each time step,
for every shell $s$:
\begin{enumerate}\setlength{\itemsep}{1pt}
\item \textbf{Imaginary-time propagation} of the amplitude,
$\partial_\tau u_s=-H_s u_s$~\eqref{eq:radial-itp}, relaxing $u_s$ toward the lowest
eigenfunction of $H_s$;
\item a \textbf{Poisson / shell-theorem solve} for the per-electron Hartree potential
$W_s$~\eqref{eq:radial-hartree};
\item a \textbf{boundary move}: each shell radius $R_s$ is shifted in the direction that
\emph{decreases the jump} of $u$ across it, driving the interface toward the
continuity~\eqref{eq:radial-bc} --- the Bernoulli free-boundary condition.
\end{enumerate}
After each step $u_s$ is renormalised~\eqref{eq:radial-bc}; at convergence the boundaries
are stationary and the per-electron energy is the Rayleigh quotient
$\varepsilon_s$~\eqref{eq:radial-E}, supplying the ionization energy
through~\eqref{eq:IE-direct}. Each of the three is a single one-dimensional grid sweep on
$r_i=ih$, so every atom converges in a fraction of a second and the whole periodic table
in about a minute on a laptop.

\section{Results: total energies, packing, and ionization across the periodic table}\label{sec:results}

We report three results. The electron packing (Section~\ref{sec:packing}) and the
all-electron total energies of the first three periods (Section~\ref{sec:TEexplicit})
come from the spherically-reduced (1D radial) solver applied to every electron. The
full main-group periodic table (Section~\ref{sec:frozencore}) is then reached by the
two-scale frozen-core decomposition of Section~\ref{sec:twoscale}: an inner core taken
from reference data~\cite{ClementiRoetti,Desclaux} plus the RealQM valence computed in
its pseudo-kernel. The
transition and inner-transition elements (the $d$- and $f$-blocks) enter only as closed
buried cores --- physically appropriate, since their electrons lie \emph{inside} the
valence and barely participate in chemistry --- and are not resolved individually;
``main group'' below means groups 1--2 and 13--18. Ionization energies
(Section~\ref{sec:ie}) are summarised last.

\subsection{Electron packing: 4+4, not 8}\label{sec:packing}

For each atom the shell packing is determined by total energy. Two findings are
robust across $Z=2$--$18$. \emph{First}, comparing shells of capacity 2, 4, 6, 8,
the \textbf{cap-4} packing matches reference total energies best everywhere: a
capacity-8 shell over-binds by $+9\%$ to $+22\%$ (worsening with $Z$), and a closed
octet is realised only as \textbf{4+4}, never as a single 8-shell --- the
geometrically expected result, since four unit charges pack as a stable tetrahedron
while eight on one shell do not. \emph{Second}, distributing each period's electrons
\emph{evenly} across its cap-4 shells (balanced filling) rather than greedily
reproduces the classic configurations 2+2+1 (B), 2+2+2 (C), 2+3+2 (N), 2+3+3 (O) and
tightens the energies (e.g.\ C: $+4.4\%$ greedy $\to -1.1\%$ balanced).

The organising principle is geometric: each electron is a localised cloud whose
\emph{size} is set by its distance from the kernel --- inner electrons compact, outer
ones diffuse --- and the shells are the close-packings (capacities $2$ and $4$) of these
clouds. Different packings of the same atom have different energies, and ---
importantly --- the \emph{unconstrained} energy minimum is \emph{not} the physical
configuration: in the 1D solver it favours the over-packed capacity-8 shell, the radial
model lacking the angular exclusion that makes 8-on-a-shell unstable (a first signal
that the closed shell is intrinsically three-dimensional). The physical packing is fixed
instead by the build-up (Aufbau) of the atom, electron by electron. In this paper we
therefore take the packing as the classic cap-4 configurations, selected by best fit to
the total energies, and defer a \emph{dynamical} Aufbau --- one that would select the
physical packing intrinsically --- to later work; RealQM's task here is to compute the
energies of the chosen packings accurately.

\subsection{All-electron total energies, periods 1--3}\label{sec:TEexplicit}

Table~\ref{tab:TE} reports total energies with the balanced cap-4 packing against
near-exact nonrelativistic reference values~\cite{Chakravorty}.

\begin{table}[t]
\centering
\caption{Total atomic energies (balanced cap-4 packing, spherically-reduced solver)
versus near-exact nonrelativistic references. Mean $|\mathrm{error}|=1.7\%$, maximum
$4.6\%$; most atoms within $2.5\%$.}
\label{tab:TE}
\begin{tabular}{llcrrr}
\toprule
$Z$ & atom & config & $TE$ RealQM (Ha) & reference (Ha) & error \\
\midrule
2  & He & 2         & $-2.845$   & $-2.904$   & $-2.0\%$ \\
3  & Li & 2+1       & $-7.434$   & $-7.478$   & $-0.6\%$ \\
4  & Be & 2+2       & $-14.630$  & $-14.667$  & $-0.3\%$ \\
5  & B  & 2+2+1     & $-24.025$  & $-24.654$  & $-2.6\%$ \\
6  & C  & 2+2+2     & $-37.445$  & $-37.845$  & $-1.1\%$ \\
7  & N  & 2+3+2     & $-55.263$  & $-54.589$  & $+1.2\%$ \\
8  & O  & 2+3+3     & $-75.389$  & $-75.067$  & $+0.4\%$ \\
9  & F  & 2+4+3     & $-102.494$ & $-99.734$  & $+2.8\%$ \\
10 & Ne & 2+4+4     & $-131.833$ & $-128.938$ & $+2.2\%$ \\
11 & Na & 2+4+4+1   & $-165.360$ & $-162.255$ & $+1.9\%$ \\
12 & Mg & 2+4+4+1+1 & $-201.704$ & $-200.053$ & $+0.8\%$ \\
13 & Al & 2+4+4+3   & $-248.671$ & $-242.346$ & $+2.6\%$ \\
14 & Si & 2+4+4+4   & $-302.531$ & $-289.359$ & $+4.6\%$ \\
15 & P  & 2+4+4+3+2 & $-347.535$ & $-341.259$ & $+1.8\%$ \\
16 & S  & 2+4+4+3+3 & $-415.040$ & $-398.110$ & $+4.3\%$ \\
17 & Cl & 2+4+4+4+3 & $-466.118$ & $-460.148$ & $+1.3\%$ \\
18 & Ar & 2+4+4+4+4 & $-528.000$ & $-527.540$ & $+0.1\%$ \\
\bottomrule
\end{tabular}
\end{table}

The mean relative error is $1.7\%$ and the maximum $4.6\%$ (Si, S); most atoms fall
within $2.5\%$. This is comparable to Hartree--Fock accuracy on total energy --- with
the opposite sign of error (RealQM over-binds the heavier atoms of the third period, where HF
under-binds through missing correlation) --- obtained from a 1D radial solve with no
basis sets, no antisymmetry, and no two-electron integrals.

\subsection{Frozen-core total energies: the floor and the variation}\label{sec:frozencore}

A total energy is meaningful only relative to a reference level. The absolute number is
dominated by the deep inner shells and conveys little on its own; chemistry and
periodicity live in the \emph{changes} of energy across groups and down periods. The
two-scale decomposition of Section~\ref{sec:twoscale} makes this floor explicit,
\begin{equation}\label{eq:frozencore}
TE(Z) \;=\; \underbrace{E_{\rm core}(Z)}_{\text{floor (reference)}} \;+\; \underbrace{E_{\rm val}(Z)}_{\text{RealQM valence}},
\end{equation}
with $E_{\rm core}$ the energy of the closed noble-gas-like core --- the floor that sets
the level --- and $E_{\rm val}$ the RealQM valence energy from the spherical
pseudo-kernel solve, the valence held outside the core radius $r_c$ by a Neumann
boundary $u'(r_c)=0$ (Section~\ref{sec:radialsolver}). For period~2 the floor is the two-electron He-like core,
known analytically as $-Z^2+\tfrac58 Z-0.158$; for periods 3--6 it is taken from
reference data. A single $r_c$ per period and block sets the level from one reference
energy --- the alkaline-earth valence energy for the $s$-block, the noble-gas valence
energy for the $p$-block --- so the remaining six atoms of each period are predictions.

\begin{table}[t]
\centering
\caption{Frozen-core total energies \eqref{eq:frozencore} for representative main-group
atoms versus reference. Every atom with $Z\ge 11$ lies within $0.5\%$, periods 4--6
within $0.05\%$ (floor-dominated). The period-2 $p$-block is the exception (text).}
\label{tab:frozencore}
\begin{tabular}{llrrr}
\toprule
period & atom & $TE$ RealQM (Ha) & reference (Ha) & error \\
\midrule
2 & Be & $-14.672$   & $-14.667$   & $+0.03\%$ \\
2 & Ne & $-129.321$  & $-128.938$  & $+0.30\%$ \\
3 & Na & $-162.286$  & $-162.255$  & $+0.02\%$ \\
3 & Ar & $-527.308$  & $-527.540$  & $-0.04\%$ \\
4 & K  & $-599.185$  & $-599.165$  & $+0.00\%$ \\
4 & Kr & $-2752.08$  & $-2752.06$  & $+0.00\%$ \\
5 & Xe & $-7232.19$  & $-7232.14$  & $+0.00\%$ \\
6 & Rn & $-21566.7$  & $-21566.7$  & $+0.00\%$ \\
\bottomrule
\end{tabular}
\end{table}

Table~\ref{tab:frozencore} reports the absolute total energies: every main-group atom
from $Z=11$ (Na) to $Z=86$ (Rn) is within $0.5\%$, periods 4--6 within $0.05\%$. This
precision is, as it must be, carried by the floor --- the valence is a small fraction
of the deep total, so once the core is referenced the absolute total follows. The
substantive content is therefore not the absolute number but $E_{\rm val}$ and its
variation across the table, which RealQM supplies and which the ionization energies of
Section~\ref{sec:ie} test directly.

Two points sharpen the picture. \emph{Period~2 is the exception}: its He floor is so
shallow that the valence is a large fraction of the total, and the absolute error rises
to $2$--$3\%$ for the lightest $p$-block atoms (B, F, Ne; C, N, O remain near $1\%$).
There the all-electron solve of Table~\ref{tab:TE} is actually more accurate, and the
two methods cross over cleanly --- explicit wins for period~2 (shallow floor), frozen
core wins decisively for periods 3+ (e.g.\ Si: $+0.3\%$ frozen versus $+4.6\%$
explicit). \emph{The transition elements are absorbed into the floor}: a period-4
$p$-block atom sits on an [Ar]$3d^{10}$ core, a period-6 one on an [Xe]$4f^{14}5d^{10}$
core. These buried shells are not resolved; they serve as the reference level, which is
exactly their physical role.

Table~\ref{tab:master} collects the full main-group result --- total energy and first
ionization energy, each with its error, for every atom from Li to Rn. It is the
periodic table in RealQM: total energies within $1\%$ for $Z\ge 11$ (the period-2
$p$-block the $2$--$6\%$ exception), and first ionization energies within $10$--$20\%$
for the $s$-block and the heavier $p$-block, with the light $p$-block under-binding (the
radial-octet ceiling of Section~\ref{sec:ie}). Total energies and ionization energies
are obtained from their own $r_c$ calibrations (valence-energy and ionization-energy
respectively), so the two error columns reflect two separate fits to one reference each
per period.

\begin{table}[p]
\centering
\footnotesize
\caption{Total energies (frozen core) and first ionization energies (pseudo-kernel) for
all main-group atoms, with errors. Totals: within $1\%$ for $Z\ge11$, period-2
$p$-block the exception. Ionization energies versus NIST~\cite{NIST-ASD}: $s$-block and
heavy $p$-block within $10$--$20\%$, light $p$-block the radial-octet limit.}
\label{tab:master}
\begin{tabular}{rl rr rrr}
\toprule
$Z$ & atom & $TE$ (Ha) & err & $\mathrm{IE}_1$ (eV) & NIST & err \\
\midrule
3  & Li & $-7.540$    & $+0.8\%$ & $5.39$  & $5.39$  & $0\%$   \\
4  & Be & $-14.672$   & $+0.0\%$ & $6.81$  & $9.32$  & $-27\%$ \\
5  & B  & $-26.217$   & $+6.3\%$ & $4.52$  & $8.30$  & $-45\%$ \\
6  & C  & $-39.756$   & $+5.1\%$ & $7.68$  & $11.26$ & $-32\%$ \\
7  & N  & $-57.742$   & $+5.8\%$ & $7.47$  & $14.53$ & $-49\%$ \\
8  & O  & $-77.697$   & $+3.5\%$ & $10.41$ & $13.62$ & $-24\%$ \\
9  & F  & $-102.888$  & $+3.2\%$ & $9.07$  & $17.42$ & $-48\%$ \\
10 & Ne & $-129.321$  & $+0.3\%$ & $11.29$ & $21.57$ & $-48\%$ \\
\midrule
11 & Na & $-162.286$  & $+0.0\%$ & $5.12$  & $5.14$  & $0\%$   \\
12 & Mg & $-200.057$  & $+0.0\%$ & $6.40$  & $7.65$  & $-16\%$ \\
13 & Al & $-243.126$  & $+0.3\%$ & $4.52$  & $5.99$  & $-24\%$ \\
14 & Si & $-290.235$  & $+0.3\%$ & $7.68$  & $8.15$  & $-6\%$  \\
15 & P  & $-342.725$  & $+0.4\%$ & $7.47$  & $10.49$ & $-29\%$ \\
16 & S  & $-399.211$  & $+0.3\%$ & $10.41$ & $10.36$ & $+1\%$  \\
17 & Cl & $-461.385$  & $+0.3\%$ & $9.07$  & $12.97$ & $-30\%$ \\
18 & Ar & $-527.308$  & $-0.0\%$ & $11.29$ & $15.76$ & $-28\%$ \\
\midrule
19 & K  & $-599.185$  & $+0.0\%$ & $4.33$  & $4.34$  & $0\%$   \\
20 & Ca & $-676.761$  & $+0.0\%$ & $5.24$  & $6.11$  & $-14\%$ \\
31 & Ga & $-1923.63$  & $+0.0\%$ & $4.52$  & $6.00$  & $-25\%$ \\
32 & Ge & $-2075.73$  & $+0.0\%$ & $7.68$  & $7.90$  & $-3\%$  \\
33 & As & $-2235.12$  & $+0.0\%$ & $7.47$  & $9.79$  & $-24\%$ \\
34 & Se & $-2400.58$  & $+0.0\%$ & $10.41$ & $9.75$  & $+7\%$  \\
35 & Br & $-2573.36$  & $+0.0\%$ & $9.07$  & $11.81$ & $-23\%$ \\
36 & Kr & $-2752.08$  & $+0.0\%$ & $11.29$ & $14.00$ & $-19\%$ \\
\midrule
37 & Rb & $-2938.37$  & $+0.0\%$ & $4.18$  & $4.18$  & $0\%$   \\
38 & Sr & $-3131.55$  & $-0.0\%$ & $5.03$  & $5.70$  & $-12\%$ \\
49 & In & $-5740.36$  & $+0.0\%$ & $4.41$  & $5.79$  & $-24\%$ \\
50 & Sn & $-6023.09$  & $+0.0\%$ & $7.12$  & $7.34$  & $-3\%$  \\
51 & Sb & $-6313.96$  & $+0.0\%$ & $7.18$  & $8.61$  & $-17\%$ \\
52 & Te & $-6612.09$  & $+0.0\%$ & $9.02$  & $9.01$  & $0\%$   \\
53 & I  & $-6917.86$  & $-0.0\%$ & $9.58$  & $10.45$ & $-8\%$  \\
54 & Xe & $-7232.19$  & $+0.0\%$ & $12.19$ & $12.13$ & $+1\%$  \\
\midrule
55 & Cs & $-7553.94$  & $+0.0\%$ & $3.91$  & $3.89$  & $+1\%$  \\
56 & Ba & $-7883.54$  & $-0.0\%$ & $4.66$  & $5.21$  & $-10\%$ \\
81 & Tl & $-18408.96$ & $-0.0\%$ & $4.21$  & $6.11$  & $-31\%$ \\
82 & Pb & $-19013.89$ & $-0.0\%$ & $6.85$  & $7.42$  & $-8\%$  \\
83 & Bi & $-19632.76$ & $+0.0\%$ & $6.57$  & $7.29$  & $-10\%$ \\
84 & Po & $-20263.87$ & $-0.0\%$ & $8.91$  & $8.42$  & $+6\%$  \\
85 & At & $-20909.04$ & $+0.0\%$ & $9.22$  & $9.32$  & $-1\%$  \\
86 & Rn & $-21566.72$ & $+0.0\%$ & $10.46$ & $10.75$ & $-3\%$  \\
\bottomrule
\end{tabular}
\end{table}

\subsection{Boundary of the spherical model: where radial nesting fails}\label{sec:dblock}

The spherically-reduced solver is valid exactly as long as the packing is a nested
sequence of radial shells. This holds for the first three periods --- the first 18
electrons, packed as $2,\ 4{+}4,\ 4{+}4$ --- and through the next two additions
(Z=19,20). It breaks at $Z\approx 21$, where the optimal packing ceases to be
radially nested: the electrons added there are \emph{not} the outermost, as shown
model-independently by the ionisation order (they are not the first removed). They
take a more compact, \emph{interpenetrating} arrangement that a model ordering shells
strictly by radius cannot represent, so the radial solver mis-binds and then
diverges. This is not a numerical accident but the periodic table's first departure
from radial nesting --- the region conventional orbital theory labels the d-block,
though RealQM invokes no angular-momentum orbitals, only packing. Interpenetrating
packings are intrinsically three-dimensional and require the explicit 3D valence
solve. The all-electron results of Section~\ref{sec:TEexplicit} are confined to this
radially-nested regime ($Z\le 18$); the frozen-core results of
Section~\ref{sec:frozencore} reach the full main group by absorbing the $d$- and
$f$-blocks into the reference core.

\subsection{Ionization energies}\label{sec:ie}

The per-electron route of Section~\ref{sec:model} gives the first ionization energy
directly as the valence-electron energy, $\mathrm{IE}=|\varepsilon_{\rm valence}|$,
with no separate ion calculation. For the two atoms whose valence is a single
outermost shell it is accurate (Table~\ref{tab:ie}): He and Li to ${\sim}1\%$,
validating the method.

\begin{table}[t]
\centering
\caption{First ionization energy as the valence-electron energy
$|\varepsilon_{\rm valence}|$ (balanced cap-4 packing), versus NIST~\cite{NIST-ASD}.}
\label{tab:ie}
\begin{tabular}{llrrr}
\toprule
$Z$ & atom & $\mathrm{IE}_1$ RealQM (eV) & NIST (eV) & error \\
\midrule
2 & He & 24.27 & 24.59 & $-1.3\%$ \\
3 & Li & 5.46  & 5.39  & $+1.3\%$ \\
\bottomrule
\end{tabular}
\end{table}

On a bare nucleus the eigenvalue degrades beyond Li because the outer shell is placed
at over-screened large radius. The pseudo-kernel of Section~\ref{sec:twoscale} removes
this: with the core replaced by its reference field and the valence held outside $r_c$
by the Neumann boundary $u'(r_c)=0$, the per-electron route extends across the
table. The ionization energy is the \emph{change} of energy on removing the outermost
electron --- the prototypical periodic observable --- and the model reproduces its two
defining trends: $\mathrm{IE}$ falls monotonically down every group and rises across
every period. Quantitatively (Table~\ref{tab:ieperiodic}) it is within $10$--$20\%$ for
the $s$-block (alkaline earths to a few percent) and for the heavier $p$-block (periods
5--6, several atoms within a few percent).

\begin{table}[t]
\centering
\caption{First ionization energy from the pseudo-kernel valence eigenvalue versus
NIST~\cite{NIST-ASD}, representative atoms. The down-group and across-period trends are
reproduced; the $s$-block and heavy $p$-block fall within $10$--$20\%$, the light
$p$-block under-binds (radial-octet ceiling, text).}
\label{tab:ieperiodic}
\begin{tabular}{llrrr}
\toprule
block & atom & RealQM (eV) & NIST (eV) & error \\
\midrule
$s$ (group 1)        & Na & 5.1  & 5.14  & $\sim\!0\%$ \\
$s$ (group 2)        & Mg & 6.4  & 7.65  & $-16\%$ \\
$s$ (group 2)        & Ba & 4.7  & 5.21  & $-10\%$ \\
$p$ heavy (period 5) & Sn & 7.1  & 7.34  & $-3\%$ \\
$p$ heavy (period 5) & Xe & 12.2 & 12.13 & $+1\%$ \\
$p$ light (period 3) & Ar & 8.0  & 15.76 & $-49\%$ \\
$p$ light (period 2) & Ne & 10.9 & 21.57 & $-50\%$ \\
\bottomrule
\end{tabular}
\end{table}

\begin{figure}[t]
\centering
\includegraphics[width=\textwidth]{IE_all_final.png}
\caption{First ionization energy across the periodic table: RealQM (coloured by block ---
$s$ blue, $p$ red, $d$ green) versus NIST~\cite{NIST-ASD} (black). The periodic law is
reproduced throughout --- the sawtooth of the main group and the flat transition-metal
plateau. The $s$-block and heavy $p$-block are quantitative ($10$--$20\%$); the light
$p$-block and the $d$-block are offset (the radial-octet ceiling and the missing $4s/3d$
penetration), with the \emph{shapes} correct.}
\label{fig:ie}
\end{figure}

The light $p$-block (periods 2--4) is the documented exception, under-binding by
$35$--$50\%$. The cause is structural and was anticipated in Section~\ref{sec:dblock}:
the closed octet, realised radially as $4{+}4$, places its outer tetrahedron at too
large a radius. The radial octet has a hard ionization ceiling near $13$~eV --- it
cannot reach Ne ($21.6$~eV) or Ar ($15.8$~eV) at any $r_c$ --- while the heavier noble
gases, whose ionization energies fall below that ceiling, are recovered, with the
crossover near Xe. This is not a tunable error but the signature of missing angular
structure: the physical octet is two tetrahedra at one radius (compact, angularly
separated), which only the 3D valence solve represents. Total energies are insensitive
to it --- the misplaced outer electron is a small fraction of the floored total
(Section~\ref{sec:frozencore}) --- but ionization energies, being precisely that
electron, are not. Consistent with the floor argument of
Section~\ref{sec:frozencore}, balanced filling (which spreads the outer electrons)
slightly worsens the eigenvalue even as it improves the total energy.

\section{Discussion: packing versus orbitals}\label{sec:discussion}

The shells $s,p,d,f$ that organise the conventional periodic table are the
\emph{eigenmodes of the one-electron hydrogen atom} --- the angular-momentum
eigenstates $R_{n\ell}(r)\,Y_{\ell m}(\theta,\phi)$ of $-\tfrac12\nabla^2-1/r$, their
capacities $2(2\ell+1)=2,6,10,14$ merely counting the spherical harmonics at each $\ell$
(times spin). This is a single-electron, non-interacting spectrum, exact for hydrogen and
for nothing else; the $\ell$-degeneracy is the special symmetry of the pure $1/r$
potential, already lifted by screening in every many-electron atom. The shells enter the
periodic table by \emph{imposition}: the orbital approximation places each electron in a
hydrogenic mode, fills them by the Aufbau order, and enforces exclusion through the
antisymmetry of a Slater determinant.

How fully does this \emph{explain} the periodic table? The rigorous content is
single-electron: the $s,p,d,f$ capacities $2(2\ell+1)$ are the spherical-harmonic counts
of the one-electron hydrogen spectrum, and the magic numbers $2n^2$ their sums. The
many-electron structure, however, rests on rules that are \emph{descriptive} rather than
derived --- the Aufbau/Madelung ($n+\ell$) filling order and Hund's rules are empirical
regularities, not consequences of the Schr\"odinger equation, and the $4s$-before-$3d$
ordering in particular has no first-principles derivation within the orbital picture. So
standard quantum mechanics \emph{reproduces} the table with great economy, but its
physical explanation is layered: exact angular momentum for one electron, plus fitted
filling rules for many. RealQM relocates the capacities to geometry --- $2$ and $4$ and
the octet $4{+}4$ follow from close-packing under Coulomb repulsion, the shell sizes from
the size-versus-distance rule --- but, as in Section~\ref{sec:packing}, the selection of
the physical packing (its own Aufbau) is likewise not yet derived and is here supplied by
fit. Both frameworks thus explain the \emph{capacities} from genuine physics (harmonic
counting; close-packing) while still leaning on an empirical element for the \emph{order}
in which many-electron configurations are built.

RealQM proceeds differently. There are no eigenmodes: the electrons are localised,
non-overlapping charge clouds, and the ground state minimises the Coulomb energy under
spatial exclusion. What emerges is geometric \emph{packing} --- the close-packing of
mutually repelling charges on a shell (the Thomson problem), whose small-number optima are
the antipodal pair (cap-2) and the tetrahedron (cap-4), the octet realised as $4{+}4$. The
essential difference is in how exclusion is realised: antisymmetry in $3N$-dimensional
configuration space (abstract) versus non-overlap of domains in three-space (physical).
The angular eigenmodes --- the $p$ dumbbells, the $d$ cloverleaves --- follow from the
former; the tetrahedra follow from the latter.

A regularity of the packing, borne out by the table, sharpens the point: \textbf{every
period closes as $4{+}4$, never as a single $4$}. The two are distinct objects. A single
tetrahedron has $T_d$ symmetry and is \emph{directional} --- four vertices, no inversion
centre --- whereas $4{+}4$ places the second tetrahedron as the \emph{dual} of the first,
the pair forming the cube ($O_h$ symmetry, the eight vertices of the \emph{stella
octangula}). A closed, inert shell must be spherically symmetric, and only the cube
achieves this; the single tetrahedron is its directional, \emph{reactive} half. The octet
rule is then geometric --- an atom closes its directional $4$ into the spherical $4{+}4$
--- and carbon's tetravalence is precisely the four bond directions of its single
tetrahedron, completed to the cube by sharing one electron with each of four partners.
(Equivalently, $8$ is two \emph{dual} tetrahedra forming one cubic shell, of which the
radial $4{+}4$ is the 1D shadow --- which is why the radial octet under-binds the
ionization energy of Section~\ref{sec:ie}.)

RealQM moreover gives its \emph{own} rationale for the capacities $2n^2$, with no appeal
to the harmonic count. Each shell is two \emph{dual half-shells} --- the pairing realised
geometrically, the two interpenetrating tetrahedra of the octet being the $n=2$ instance
--- and each half divides regularly into $n^2$ cells as nested rings of
$1,3,5,\dots,2n{-}1$ points, so $n^2=\sum_{\ell=0}^{n-1}(2\ell+1)$ and the shell holds
$2n^2$. The ring counts $2\ell+1$ are the geometric counterpart of the angular
degeneracy: for $n=2$ a half is $1{+}3$ --- an apex over a triangle, i.e.\ a tetrahedron
--- and the two duals form the cube. Packing alone thus yields both the magic numbers and
their $2\ell+1$ substructure, with no angular-momentum operator.

This also explains why $8$ is privileged. Only $n=2$ gives a half ($1{+}3=4$) whose
double is a \emph{Platonic} closed surface (the cube); for $n\ge3$ the half
($9,16,\dots$) is a regular division with no unique polyhedron. So $8$ is a closed,
spherically-symmetric \emph{surface}, whereas $18$ and $32$ are \emph{nested composites}
--- a cube wrapping interior $d$- and $f$-packings. Three facts follow at once: every
period closes as the octet $4{+}4$ (the only available closed surface), the $d$- and
$f$-blocks are \emph{buried} (no surface carries them), and the long periods are octet
surfaces over interior shells, not larger noble-gas surfaces.

The two pictures agree on the magic numbers $2,8,18,32$ because both count arrangements on
nested spheres in three dimensions --- the same geometry read spectrally (the harmonic
count $2n^2$) or geometrically (the packing capacity) --- but they decompose them
differently: $2{+}6{+}10{+}14$ into sub-shells versus $4{+}4$, $4{+}4{+}2$,
$4{+}4{+}4{+}2$ into tetrahedra and pairs (the buried $d$- and $f$-shells of
Section~\ref{sec:frozencore}). A closed shell is spherically symmetric either way; the
pictures part company in the angular correlation of open shells and, decisively, in
\emph{bonding geometry}. There the geometric picture is the one chemistry already uses:
VSEPR and $sp^3$ hybridisation are geometric repairs grafted onto the orbital model
precisely because the bare eigenmodes do not give molecular shape --- one must
\emph{hybridise} $s$ and $p$ into four tetrahedral lobes to recover what the cap-4 packing
yields directly. RealQM does not claim the orbitals are wrong; it locates them as the
\emph{spectral shadow of the one-electron problem}, while the structure of the real
interacting atom --- and of the chemistry built on it --- is geometric packing.

\section{Numerical properties of the solver}\label{sec:convergence}

\paragraph{Spatial discretisation.} We measure the total energy under grid refinement
at fixed box size, estimating the observed order $p$ from $|E_h-E_{h/2}|\sim h^p$
(Table~\ref{tab:hconv}). For He the convergence is clean and \textbf{first order}
($p=0.94$): the nuclear cusp --- a kink in $u$ at $r=0$ --- caps the formally
second-order stencil at $\mathcal{O}(h)$, which is also why the light-atom total
energies retain a ${\sim}1$--$2\%$ residual at the production grid.

\begin{table}[t]
\centering
\caption{Grid-refinement convergence of the total energy at fixed box (He). The order is
first ($p\approx 1$), cusp-limited. Heavier atoms require the grid to resolve the inner
$1s$ ($r\sim 1/Z$): at this coarse $N$ the Ne and Ar cores are under-resolved and the
energies are meaningless, so the production runs use $N=400$ (see text).}
\label{tab:hconv}
\begin{tabular}{lccccc}
\toprule
Atom & $N$ & $h$ (a.u.) & $E_{\rm total}$ (Ha) & $|E_N-E_{2N}|$ & order $p$ \\
\midrule
He & 64  & 0.078 & $-2.79543$ & ---     & --- \\
He & 128 & 0.039 & $-2.82554$ & 0.03012 & --- \\
He & 256 & 0.020 & $-2.84121$ & 0.01566 & $0.94$ \\
\bottomrule
\end{tabular}
\end{table}

The inner $1s$ contracts as $1/Z$, so a uniform grid must be refined in proportion to
resolve it: at the coarse $N$ above the Ne ($1s$ at $r\!\approx\!0.1$) and Ar
($r\!\approx\!0.056$) cores are unresolved and the energies diverge. This $1/Z$
resolution requirement is the numerical face of the radial-nesting boundary
(Section~\ref{sec:dblock}) and the practical reason heavy-atom total energies are
obtained through the frozen core (Section~\ref{sec:frozencore}) rather than by direct
all-electron refinement.

\paragraph{Iteration and free-boundary convergence.} The imaginary-time residual
$\|\dot u_s\|$ and the boundary jumps $|u_s(R_s)-u_{s+1}(R_s)|$ both decay to steady
state; at convergence the jumps vanish, verifying the Bernoulli (continuity) condition at
every interface.

\paragraph{Throughput.} Each step is an $\mathcal{O}(N)$ one-dimensional grid sweep over
the radial mesh; a single atom converges in a fraction of a second and the full
main-group table in about a minute on a laptop CPU --- no basis sets, no two-electron
integrals, no $N$-body configuration space.

\section{Limitations}\label{sec:limitations}

Only the stationary (converged) solution is used, so the temporal order of the
imaginary-time relaxation is immaterial. The relevant limitation is the \emph{spatial}
discretisation: it is low order and cusp-limited (first order at the nucleus,
Section~\ref{sec:convergence}), so the light-atom total energies retain a ${\sim}1$--$2\%$
residual that a higher-order stencil or a logarithmic radial mesh would reduce. The core/valence split is a
modelling choice. For the transition metals the valence is configured as the $d$-shell
nested \emph{inside} the $s$-shell, which reproduces the correct ionization ordering
($4s$ removed before $3d$) and, with a reference core, total energies to ${\sim}1\%$;
the quantitative $d$-block ionization \emph{trends}, which require $4s$/$3d$
interpenetration, remain a 3D effect. Open-shell and strongly non-spherical
configurations test the angular representation of the 3D valence solve and are flagged
where relevant. Electron affinities (negative ions) are a known difficult case and are
outside the scope of this paper.

\section{Conclusion}\label{sec:conclusion}

The results organise into a single statement about symmetry. The \emph{inner} shells of
every atom --- the closed, buried cap-2/cap-4 shells that carry most of the electrons
and nearly all of the energy --- are always spherically symmetric, and the
spherically-reduced solver computes them exactly. \emph{Most valence shells are
spherically symmetric too}: the $s$-block, the single outermost electron, the
radially-nested fillings, and the ionization \emph{ordering} throughout (including $4s$
before $3d$) are all reproduced in spherical symmetry, and with a reference core the
total energies follow to ${\sim}1\%$ across the whole periodic table, transition metals
included. Only \emph{two specific valence geometries break spherical symmetry} and
require the full 3D solve for their quantitative ionization energy: the \textbf{closed
octet} --- two tetrahedra at one radius, which the radial $4{+}4$ under-binds against a
${\sim}13$~eV ceiling (the light $p$-block) --- and the \textbf{interpenetrating
$4s/3d$} valence of the transition metals, where the outer $s$ must penetrate the inner
$d$ to give the rising ionization trend. Both are configurations in which electrons
share one radius and are distinguished only by angle, which a radially-nested model
cannot place.

The architecture and energetics of the periodic table are therefore \emph{radial}, and
are captured by RealQM in spherical symmetry; its finest valence energetics are
\emph{angular}, and are the task of the full 3D solve of Ref.~\cite{Johnson-RealQM}.
The spherical model is not a provisional approximation awaiting the 3D one --- it is the
correct description wherever the physics is radial, which is almost everywhere.

\section{Code and reproducibility}\label{sec:code}

The complete implementation, the validation suite, and an interactive gallery are
publicly available:
\begin{itemize}\setlength{\itemsep}{1pt}
  \item Source code: \url{https://github.com/Claes542/RealMolecule}
  \item Interactive gallery: \url{https://claes542.github.io/RealMolecule/}
\end{itemize}
Every figure and table is reproducible from the published source on the hardware
stated in Section~\ref{sec:convergence}.

\begin{thebibliography}{9}
\bibitem{Johnson-RealQM} C.~Johnson, \emph{Real Quantum Mechanics}, submitted to
\emph{Annales de la Fondation Louis de Broglie} (2026); preprint and code at the
repository of Section~\ref{sec:code}.
\bibitem{NIST-ASD} A.~Kramida, Yu.~Ralchenko, J.~Reader, and NIST ASD Team,
\emph{NIST Atomic Spectra Database} (ver.\ 5), National Institute of Standards and
Technology, Gaithersburg, MD.
\bibitem{KS65} W.~Kohn and L.~J.~Sham, \emph{Self-consistent equations including
exchange and correlation effects}, Phys.\ Rev.\ \textbf{140} (1965) A1133.
\bibitem{ClementiRoetti} E.~Clementi and C.~Roetti, \emph{Roothaan--Hartree--Fock
atomic wavefunctions: basis functions and their coefficients for ground and certain
excited states of neutral and ionized atoms, $Z\le54$}, At.\ Data Nucl.\ Data Tables
\textbf{14} (1974) 177--478.
\bibitem{Desclaux} J.~P.~Desclaux, \emph{Relativistic Dirac--Fock expectation values
for atoms with $Z=1$ to $Z=120$}, At.\ Data Nucl.\ Data Tables \textbf{12} (1973)
311--406.
\bibitem{Chakravorty} S.~J.~Chakravorty, S.~R.~Gwaltney, E.~R.~Davidson,
F.~A.~Parpia, and C.~Froese~Fischer, \emph{Ground-state correlation energies for
atomic ions with 3 to 18 electrons}, Phys.\ Rev.\ A \textbf{47} (1993) 3649.
\bibitem{HK64} P.~Hohenberg and W.~Kohn, \emph{Inhomogeneous electron gas},
Phys.\ Rev.\ \textbf{136} (1964) B864.
\bibitem{SzaboOstlund} A.~Szabo and N.~S.~Ostlund, \emph{Modern Quantum Chemistry:
Introduction to Advanced Electronic Structure Theory}, Dover, New York (1996).
\bibitem{Beck2000} T.~L.~Beck, \emph{Real-space mesh techniques in density-functional
theory}, Rev.\ Mod.\ Phys.\ \textbf{72} (2000) 1041.
\end{thebibliography}

\end{document}
